
\documentclass{article}
\usepackage{colortbl}
\usepackage{makecell}
\usepackage{multirow}
\usepackage{supertabular}

\begin{document}

\newcounter{utterance}

\centering \large Interaction Transcript for game `dond', experiment `coop\_en', episode 11 with Qwen3{-}32B{-}t0.0.
\vspace{24pt}

{ \footnotesize  \setcounter{utterance}{1}
\setlength{\tabcolsep}{0pt}
\begin{supertabular}{c@{$\;$}|p{.15\linewidth}@{}p{.15\linewidth}p{.15\linewidth}p{.15\linewidth}p{.15\linewidth}p{.15\linewidth}}
    \# & $\;$A & \multicolumn{4}{c}{Game Master} & $\;\:$B\\
    \hline

    \theutterance \stepcounter{utterance}  
    & & \multicolumn{4}{p{0.6\linewidth}}{
        \cellcolor[rgb]{0.9,0.9,0.9}{
            \makecell[{{p{\linewidth}}}]{
                \texttt{\tiny{[P1$\langle$GM]}}
                \texttt{You are playing a cooperative negotiation game in which you have to agree on how to divide a set of items among you and another player.} \\
\\ 
\texttt{Rules:} \\
\texttt{(a) You and the other player are given a set of items. Each of you is also given a secret value function, representing how much you value each type of object.} \\
\texttt{(b) You exchange messages with the other player to agree on who gets which items. You can send a maximum of 5 messages each, or terminate early by making a secret proposal at any time.} \\
\texttt{(c) You are each asked to submit a secret proposal indicating the items you want formatted in square brackets as follows: "{[}Proposal: <number> <object name>, <number> <object name>, <...>{]}"} \\
\texttt{(d) If your proposals are complementary, i.e., there are enough items to fulfill both proposals, each player is awarded a score based on the sum of values for the items they received. Otherwise, both of you get zero points.} \\
\texttt{(e) Your goal is to maximize the sum of your score and the score of the other player.} \\
\\ 
\texttt{Let us start.} \\
\\ 
\texttt{The set of available items is:} \\
\\ 
\texttt{1 ring, 3 fruits, 1 chalk} \\
\\ 
\texttt{Your secret value function is:} \\
\\ 
\texttt{ring: 6, fruit: 1, chalk: 1} \\
\\ 
\texttt{IMPORTANT: Your messages, unless it is the secret proposal, are directly transmitted to the other player, so do not include any response to the rules or text announcing your message. To make a secret proposal, use the indicated format. Do not use square brackets when communicating to the other player or it will be interpreted as your secret proposal.} \\
            }
        }
    }
    & \\ \\

    \theutterance \stepcounter{utterance}  
    & \multicolumn{4}{p{0.6\linewidth}}{
        \cellcolor[rgb]{0.9,0.9,0.9}{
            \makecell[{{p{\linewidth}}}]{
                \texttt{\tiny{[P1$\rangle$GM]}}
                \texttt{<think>} \\
\texttt{Okay, let's see. I need to divide the items between me and another player. The items are 1 ring, 3 fruits, and 1 chalk. My value function is ring: 6, fruit: 1, chalk: 1. The other player's value function is secret, so I don't know what they value. My goal is to maximize the total score (mine plus theirs) if our proposals are complementary.} \\
\\ 
\texttt{First, I should figure out what the other player might value. Since I don't have their value function, I need to make some assumptions or ask for information. But since I can only send 5 messages, I need to be efficient. Maybe I can ask them what they value most. But wait, the messages are directly transmitted, so I can't include any meta instructions. So I need to phrase it as a message.} \\
\\ 
\texttt{Wait, the rules say that messages are exchanged, and I can send up to 5 messages each. But the other player is also sending messages. So maybe I should start by asking them about their preferences. For example, I could say, "What do you value most?" But since the other player's value function is secret, they might not tell me directly. Alternatively, maybe I can propose a division and see if they accept or counteroffer.} \\
\\ 
\texttt{But the key is that we need to make secret proposals that are complementary. If our proposals don't conflict, we get our scores. Otherwise, zero. So I need to find a way to divide the items so that both of our desired items don't overlap too much.} \\
\\ 
\texttt{Looking at the items: there's 1 ring, 3 fruits, 1 chalk. My highest value is the ring (6), then fruits and chalk are 1 each. So I probably want the ring. But if the other player also wants the ring, that's a conflict. But maybe they value fruits more? Or chalk?} \\
\\ 
\texttt{If I take the ring, then the other player can take the fruits and chalk. But there are 3 fruits. If the other player values fruits highly, maybe they want all 3. But if they value chalk more, maybe they want that. Alternatively, maybe they want a combination.} \\
\\ 
\texttt{Alternatively, maybe the other player's value function is such that they value the ring more than me. But since I don't know, I need to negotiate.} \\
\\ 
\texttt{Let me think of possible scenarios. Suppose the other player values the ring as} \\
            }
        }
    }
    & & \\ \\

    \theutterance \stepcounter{utterance}  
    & & & \multicolumn{2}{p{0.3\linewidth}}{
        \cellcolor[rgb]{0.9,0.9,0.9}{
            \makecell[{{p{\linewidth}}}]{
                \texttt{\tiny{[GM$|$GM]}}
                \texttt{continue} \\
            }
        }
    }
    & & \\ \\

    \theutterance \stepcounter{utterance}  
    & & & \multicolumn{2}{p{0.3\linewidth}}{
        \cellcolor[rgb]{0.9,0.9,0.9}{
            \makecell[{{p{\linewidth}}}]{
                \texttt{\tiny{[GM$|$GM]}}
                \texttt{<think>} \\
\texttt{Okay, let's see. I need to divide the items between me and another player. The items are 1 ring, 3 fruits, and 1 chalk. My value function is ring: 6, fruit: 1, chalk: 1. The other player's value function is secret, so I don't know what they value. My goal is to maximize the total score (mine plus theirs) if our proposals are complementary.} \\
\\ 
\texttt{First, I should figure out what the other player might value. Since I don't have their value function, I need to make some assumptions or ask for information. But since I can only send 5 messages, I need to be efficient. Maybe I can ask them what they value most. But wait, the messages are directly transmitted, so I can't include any meta instructions. So I need to phrase it as a message.} \\
\\ 
\texttt{Wait, the rules say that messages are exchanged, and I can send up to 5 messages each. But the other player is also sending messages. So maybe I should start by asking them about their preferences. For example, I could say, "What do you value most?" But since the other player's value function is secret, they might not tell me directly. Alternatively, maybe I can propose a division and see if they accept or counteroffer.} \\
\\ 
\texttt{But the key is that we need to make secret proposals that are complementary. If our proposals don't conflict, we get our scores. Otherwise, zero. So I need to find a way to divide the items so that both of our desired items don't overlap too much.} \\
\\ 
\texttt{Looking at the items: there's 1 ring, 3 fruits, 1 chalk. My highest value is the ring (6), then fruits and chalk are 1 each. So I probably want the ring. But if the other player also wants the ring, that's a conflict. But maybe they value fruits more? Or chalk?} \\
\\ 
\texttt{If I take the ring, then the other player can take the fruits and chalk. But there are 3 fruits. If the other player values fruits highly, maybe they want all 3. But if they value chalk more, maybe they want that. Alternatively, maybe they want a combination.} \\
\\ 
\texttt{Alternatively, maybe the other player's value function is such that they value the ring more than me. But since I don't know, I need to negotiate.} \\
\\ 
\texttt{Let me think of possible scenarios. Suppose the other player values the ring as} \\
            }
        }
    }
    & & \\ \\

    \theutterance \stepcounter{utterance}  
    & & \multicolumn{4}{p{0.6\linewidth}}{
        \cellcolor[rgb]{0.9,0.9,0.9}{
            \makecell[{{p{\linewidth}}}]{
                \texttt{\tiny{[GM$\rangle$P2]}}
                \texttt{You are playing a cooperative negotiation game in which you have to agree on how to divide a set of items among you and another player.} \\
\\ 
\texttt{Rules:} \\
\texttt{(a) You and the other player are given a set of items. Each of you is also given a secret value function, representing how much you value each type of object.} \\
\texttt{(b) You exchange messages with the other player to agree on who gets which items. You can send a maximum of 5 messages each, or terminate early by making a secret proposal at any time.} \\
\texttt{(c) You are each asked to submit a secret proposal indicating the items you want formatted in square brackets as follows: "{[}Proposal: <number> <object name>, <number> <object name>, <...>{]}"} \\
\texttt{(d) If your proposals are complementary, i.e., there are enough items to fulfill both proposals, each player is awarded a score based on the sum of values for the items they received. Otherwise, both of you get zero points.} \\
\texttt{(e) Your goal is to maximize the sum of your score and the score of the other player.} \\
\\ 
\texttt{Let us start.} \\
\\ 
\texttt{The set of available items is:} \\
\\ 
\texttt{1 ring, 3 fruits, 1 chalk} \\
\\ 
\texttt{Your secret value function is:} \\
\\ 
\texttt{ring: 6, fruit: 1, chalk: 1} \\
\\ 
\texttt{IMPORTANT: Your messages, unless it is the secret proposal, are directly transmitted to the other player, so do not include any response to the rules or text announcing your message. To make a secret proposal, use the indicated format. Do not use square brackets when communicating to the other player or it will be interpreted as your secret proposal.} \\
\\ 
\\ 
\texttt{<think>} \\
\texttt{Okay, let's see. I need to divide the items between me and another player. The items are 1 ring, 3 fruits, and 1 chalk. My value function is ring: 6, fruit: 1, chalk: 1. The other player's value function is secret, so I don't know what they value. My goal is to maximize the total score (mine plus theirs) if our proposals are complementary.} \\
\\ 
\texttt{First, I should figure out what the other player might value. Since I don't have their value function, I need to make some assumptions or ask for information. But since I can only send 5 messages, I need to be efficient. Maybe I can ask them what they value most. But wait, the messages are directly transmitted, so I can't include any meta instructions. So I need to phrase it as a message.} \\
\\ 
\texttt{Wait, the rules say that messages are exchanged, and I can send up to 5 messages each. But the other player is also sending messages. So maybe I should start by asking them about their preferences. For example, I could say, "What do you value most?" But since the other player's value function is secret, they might not tell me directly. Alternatively, maybe I can propose a division and see if they accept or counteroffer.} \\
\\ 
\texttt{But the key is that we need to make secret proposals that are complementary. If our proposals don't conflict, we get our scores. Otherwise, zero. So I need to find a way to divide the items so that both of our desired items don't overlap too much.} \\
\\ 
\texttt{Looking at the items: there's 1 ring, 3 fruits, 1 chalk. My highest value is the ring (6), then fruits and chalk are 1 each. So I probably want the ring. But if the other player also wants the ring, that's a conflict. But maybe they value fruits more? Or chalk?} \\
\\ 
\texttt{If I take the ring, then the other player can take the fruits and chalk. But there are 3 fruits. If the other player values fruits highly, maybe they want all 3. But if they value chalk more, maybe they want that. Alternatively, maybe they want a combination.} \\
\\ 
\texttt{Alternatively, maybe the other player's value function is such that they value the ring more than me. But since I don't know, I need to negotiate.} \\
\\ 
\texttt{Let me think of possible scenarios. Suppose the other player values the ring as} \\
            }
        }
    }
    & \\ \\

    \theutterance \stepcounter{utterance}  
    & & & \multicolumn{4}{p{0.6\linewidth}}{
        \cellcolor[rgb]{0.9,0.9,0.9}{
            \makecell[{{p{\linewidth}}}]{
                \texttt{\tiny{[GM$\langle$P2]}}
                \texttt{<think>} \\
\texttt{Okay, let's try to figure out the best way to approach this. The items are 1 ring, 3 fruits, and 1 chalk. My values are ring:6, fruit:1, chalk:1. The other player's values are unknown. The goal is to maximize the total score (mine plus theirs) if our proposals are complementary.} \\
\\ 
\texttt{First, I need to consider possible value functions the other player might have. Since the items are ring, fruits, and chalk, common valuations could be:} \\
\\ 
\texttt{1. They might value the ring highly (like me), maybe even more.} \\
\texttt{2. They might value fruits more, especially since there are 3 of them.} \\
\texttt{3. They might value chalk for some reason, maybe it's rare or has a special use.} \\
\\ 
\texttt{Since I don't know their values, I need to find a division that's efficient regardless. The key is to maximize the total value, so I should aim for a division where the sum of both players' values is highest.} \\
\\ 
\texttt{If I take the ring (6 for me), then the remaining items are 3 fruits and 1 chalk. If the other player values fruits at, say, 2 each, then giving them all 3 fruits would give them 6, and chalk maybe 1. But if they value fruits lower, maybe 1 each, then 3 fruits would be 3. But if they value chalk higher, maybe they prefer that.} \\
\\ 
\texttt{Alternatively, if the other player also wants the ring, there's a conflict. But since there's only one ring, we can't both have it. So maybe I should propose taking the ring and offer the other player the fruits and chalk. But if they need the ring, that's a problem.} \\
\\ 
\texttt{Wait, but the secret proposal needs to be complementary. So if I propose taking the ring, and they propose taking the fruits and chalk, that's okay. But if they also propose taking the ring, then our proposals conflict, and we both get zero.} \\
\\ 
\texttt{So maybe I should try to find out what they value. But since I can't ask directly (since messages are sent as is), maybe I can make a proposal that allows for different possibilities.} \\
\\ 
\texttt{Alternatively, maybe I can make a proposal that gives me the ring and the other player the rest. Let's see: my value would be 6 (ring) + 0 (since I don't take fruits or chalk). The other player's value would depend} \\
            }
        }
    }
     \\ \\

    \theutterance \stepcounter{utterance}  
    & & & \multicolumn{2}{p{0.3\linewidth}}{
        \cellcolor[rgb]{0.9,0.9,0.9}{
            \makecell[{{p{\linewidth}}}]{
                \texttt{\tiny{[GM$|$GM]}}
                \texttt{continue} \\
            }
        }
    }
    & & \\ \\

    \theutterance \stepcounter{utterance}  
    & & & \multicolumn{2}{p{0.3\linewidth}}{
        \cellcolor[rgb]{0.9,0.9,0.9}{
            \makecell[{{p{\linewidth}}}]{
                \texttt{\tiny{[GM$|$GM]}}
                \texttt{<think>} \\
\texttt{Okay, let's try to figure out the best way to approach this. The items are 1 ring, 3 fruits, and 1 chalk. My values are ring:6, fruit:1, chalk:1. The other player's values are unknown. The goal is to maximize the total score (mine plus theirs) if our proposals are complementary.} \\
\\ 
\texttt{First, I need to consider possible value functions the other player might have. Since the items are ring, fruits, and chalk, common valuations could be:} \\
\\ 
\texttt{1. They might value the ring highly (like me), maybe even more.} \\
\texttt{2. They might value fruits more, especially since there are 3 of them.} \\
\texttt{3. They might value chalk for some reason, maybe it's rare or has a special use.} \\
\\ 
\texttt{Since I don't know their values, I need to find a division that's efficient regardless. The key is to maximize the total value, so I should aim for a division where the sum of both players' values is highest.} \\
\\ 
\texttt{If I take the ring (6 for me), then the remaining items are 3 fruits and 1 chalk. If the other player values fruits at, say, 2 each, then giving them all 3 fruits would give them 6, and chalk maybe 1. But if they value fruits lower, maybe 1 each, then 3 fruits would be 3. But if they value chalk higher, maybe they prefer that.} \\
\\ 
\texttt{Alternatively, if the other player also wants the ring, there's a conflict. But since there's only one ring, we can't both have it. So maybe I should propose taking the ring and offer the other player the fruits and chalk. But if they need the ring, that's a problem.} \\
\\ 
\texttt{Wait, but the secret proposal needs to be complementary. So if I propose taking the ring, and they propose taking the fruits and chalk, that's okay. But if they also propose taking the ring, then our proposals conflict, and we both get zero.} \\
\\ 
\texttt{So maybe I should try to find out what they value. But since I can't ask directly (since messages are sent as is), maybe I can make a proposal that allows for different possibilities.} \\
\\ 
\texttt{Alternatively, maybe I can make a proposal that gives me the ring and the other player the rest. Let's see: my value would be 6 (ring) + 0 (since I don't take fruits or chalk). The other player's value would depend} \\
            }
        }
    }
    & & \\ \\

    \theutterance \stepcounter{utterance}  
    & & \multicolumn{4}{p{0.6\linewidth}}{
        \cellcolor[rgb]{0.9,0.9,0.9}{
            \makecell[{{p{\linewidth}}}]{
                \texttt{\tiny{[P1$\langle$GM]}}
                \texttt{<think>} \\
\texttt{Okay, let's try to figure out the best way to approach this. The items are 1 ring, 3 fruits, and 1 chalk. My values are ring:6, fruit:1, chalk:1. The other player's values are unknown. The goal is to maximize the total score (mine plus theirs) if our proposals are complementary.} \\
\\ 
\texttt{First, I need to consider possible value functions the other player might have. Since the items are ring, fruits, and chalk, common valuations could be:} \\
\\ 
\texttt{1. They might value the ring highly (like me), maybe even more.} \\
\texttt{2. They might value fruits more, especially since there are 3 of them.} \\
\texttt{3. They might value chalk for some reason, maybe it's rare or has a special use.} \\
\\ 
\texttt{Since I don't know their values, I need to find a division that's efficient regardless. The key is to maximize the total value, so I should aim for a division where the sum of both players' values is highest.} \\
\\ 
\texttt{If I take the ring (6 for me), then the remaining items are 3 fruits and 1 chalk. If the other player values fruits at, say, 2 each, then giving them all 3 fruits would give them 6, and chalk maybe 1. But if they value fruits lower, maybe 1 each, then 3 fruits would be 3. But if they value chalk higher, maybe they prefer that.} \\
\\ 
\texttt{Alternatively, if the other player also wants the ring, there's a conflict. But since there's only one ring, we can't both have it. So maybe I should propose taking the ring and offer the other player the fruits and chalk. But if they need the ring, that's a problem.} \\
\\ 
\texttt{Wait, but the secret proposal needs to be complementary. So if I propose taking the ring, and they propose taking the fruits and chalk, that's okay. But if they also propose taking the ring, then our proposals conflict, and we both get zero.} \\
\\ 
\texttt{So maybe I should try to find out what they value. But since I can't ask directly (since messages are sent as is), maybe I can make a proposal that allows for different possibilities.} \\
\\ 
\texttt{Alternatively, maybe I can make a proposal that gives me the ring and the other player the rest. Let's see: my value would be 6 (ring) + 0 (since I don't take fruits or chalk). The other player's value would depend} \\
            }
        }
    }
    & \\ \\

    \theutterance \stepcounter{utterance}  
    & \multicolumn{4}{p{0.6\linewidth}}{
        \cellcolor[rgb]{0.9,0.9,0.9}{
            \makecell[{{p{\linewidth}}}]{
                \texttt{\tiny{[P1$\rangle$GM]}}
                \texttt{<think>} \\
\texttt{Okay, let's try to figure out the best way to approach this. The items are 1 ring, 3 fruits, and 1 chalk. My values are ring:6, fruit:1, chalk:1. The other player's values are unknown. The goal is to maximize the total score (mine plus theirs) if our proposals are complementary.} \\
\\ 
\texttt{First, I need to consider possible value functions the other player might have. Since the items are ring, fruits, and chalk, common valuations could be:} \\
\\ 
\texttt{1. They might value the ring highly (like me), maybe even more.} \\
\texttt{2. They might value fruits more, especially since there are 3 of them.} \\
\texttt{3. They might value chalk for some reason, maybe it's rare or has a special use.} \\
\\ 
\texttt{Since I don't know their values, I need to find a division that's efficient regardless. The key is to maximize the total value, so I should aim for a division where the sum of both players' values is highest.} \\
\\ 
\texttt{If I take the ring (6 for me), then the remaining items are 3 fruits and 1 chalk. If the other player values fruits at, say, 2 each, then giving them all 3 fruits would give them 6, and chalk maybe 1. But if they value fruits lower, maybe 1 each, then 3 fruits would be 3. But if they value chalk higher, maybe they prefer that.} \\
\\ 
\texttt{Alternatively, if the other player also wants the ring, there's a conflict. But since there's only one ring, we can't both have it. So maybe I should propose taking the ring and offer the other player the fruits and chalk. But if they need the ring, that's a problem.} \\
\\ 
\texttt{Wait, but the secret proposal needs to be complementary. So if I propose taking the ring, and they propose taking the fruits and chalk, that's okay. But if they also propose taking the ring, then our proposals conflict, and we both get zero.} \\
\\ 
\texttt{So maybe I should try to find out what they value. But since I can't ask directly (since messages are sent as is), maybe I can make a proposal that allows for different possibilities.} \\
\\ 
\texttt{Alternatively, maybe I can make a proposal that gives me the ring and the other player the rest. Let's see: my value would be 6 (ring) + 0 (since I don't take fruits or chalk). The other player's value would depend} \\
            }
        }
    }
    & & \\ \\

    \theutterance \stepcounter{utterance}  
    & & & \multicolumn{2}{p{0.3\linewidth}}{
        \cellcolor[rgb]{0.9,0.9,0.9}{
            \makecell[{{p{\linewidth}}}]{
                \texttt{\tiny{[GM$|$GM]}}
                \texttt{continue} \\
            }
        }
    }
    & & \\ \\

    \theutterance \stepcounter{utterance}  
    & & & \multicolumn{2}{p{0.3\linewidth}}{
        \cellcolor[rgb]{0.9,0.9,0.9}{
            \makecell[{{p{\linewidth}}}]{
                \texttt{\tiny{[GM$|$GM]}}
                \texttt{<think>} \\
\texttt{Okay, let's try to figure out the best way to approach this. The items are 1 ring, 3 fruits, and 1 chalk. My values are ring:6, fruit:1, chalk:1. The other player's values are unknown. The goal is to maximize the total score (mine plus theirs) if our proposals are complementary.} \\
\\ 
\texttt{First, I need to consider possible value functions the other player might have. Since the items are ring, fruits, and chalk, common valuations could be:} \\
\\ 
\texttt{1. They might value the ring highly (like me), maybe even more.} \\
\texttt{2. They might value fruits more, especially since there are 3 of them.} \\
\texttt{3. They might value chalk for some reason, maybe it's rare or has a special use.} \\
\\ 
\texttt{Since I don't know their values, I need to find a division that's efficient regardless. The key is to maximize the total value, so I should aim for a division where the sum of both players' values is highest.} \\
\\ 
\texttt{If I take the ring (6 for me), then the remaining items are 3 fruits and 1 chalk. If the other player values fruits at, say, 2 each, then giving them all 3 fruits would give them 6, and chalk maybe 1. But if they value fruits lower, maybe 1 each, then 3 fruits would be 3. But if they value chalk higher, maybe they prefer that.} \\
\\ 
\texttt{Alternatively, if the other player also wants the ring, there's a conflict. But since there's only one ring, we can't both have it. So maybe I should propose taking the ring and offer the other player the fruits and chalk. But if they need the ring, that's a problem.} \\
\\ 
\texttt{Wait, but the secret proposal needs to be complementary. So if I propose taking the ring, and they propose taking the fruits and chalk, that's okay. But if they also propose taking the ring, then our proposals conflict, and we both get zero.} \\
\\ 
\texttt{So maybe I should try to find out what they value. But since I can't ask directly (since messages are sent as is), maybe I can make a proposal that allows for different possibilities.} \\
\\ 
\texttt{Alternatively, maybe I can make a proposal that gives me the ring and the other player the rest. Let's see: my value would be 6 (ring) + 0 (since I don't take fruits or chalk). The other player's value would depend} \\
            }
        }
    }
    & & \\ \\

    \theutterance \stepcounter{utterance}  
    & & \multicolumn{4}{p{0.6\linewidth}}{
        \cellcolor[rgb]{0.9,0.9,0.9}{
            \makecell[{{p{\linewidth}}}]{
                \texttt{\tiny{[GM$\rangle$P2]}}
                \texttt{<think>} \\
\texttt{Okay, let's try to figure out the best way to approach this. The items are 1 ring, 3 fruits, and 1 chalk. My values are ring:6, fruit:1, chalk:1. The other player's values are unknown. The goal is to maximize the total score (mine plus theirs) if our proposals are complementary.} \\
\\ 
\texttt{First, I need to consider possible value functions the other player might have. Since the items are ring, fruits, and chalk, common valuations could be:} \\
\\ 
\texttt{1. They might value the ring highly (like me), maybe even more.} \\
\texttt{2. They might value fruits more, especially since there are 3 of them.} \\
\texttt{3. They might value chalk for some reason, maybe it's rare or has a special use.} \\
\\ 
\texttt{Since I don't know their values, I need to find a division that's efficient regardless. The key is to maximize the total value, so I should aim for a division where the sum of both players' values is highest.} \\
\\ 
\texttt{If I take the ring (6 for me), then the remaining items are 3 fruits and 1 chalk. If the other player values fruits at, say, 2 each, then giving them all 3 fruits would give them 6, and chalk maybe 1. But if they value fruits lower, maybe 1 each, then 3 fruits would be 3. But if they value chalk higher, maybe they prefer that.} \\
\\ 
\texttt{Alternatively, if the other player also wants the ring, there's a conflict. But since there's only one ring, we can't both have it. So maybe I should propose taking the ring and offer the other player the fruits and chalk. But if they need the ring, that's a problem.} \\
\\ 
\texttt{Wait, but the secret proposal needs to be complementary. So if I propose taking the ring, and they propose taking the fruits and chalk, that's okay. But if they also propose taking the ring, then our proposals conflict, and we both get zero.} \\
\\ 
\texttt{So maybe I should try to find out what they value. But since I can't ask directly (since messages are sent as is), maybe I can make a proposal that allows for different possibilities.} \\
\\ 
\texttt{Alternatively, maybe I can make a proposal that gives me the ring and the other player the rest. Let's see: my value would be 6 (ring) + 0 (since I don't take fruits or chalk). The other player's value would depend} \\
            }
        }
    }
    & \\ \\

    \theutterance \stepcounter{utterance}  
    & & & \multicolumn{4}{p{0.6\linewidth}}{
        \cellcolor[rgb]{0.9,0.9,0.9}{
            \makecell[{{p{\linewidth}}}]{
                \texttt{\tiny{[GM$\langle$P2]}}
                \texttt{<think>} \\
\texttt{Okay, let's try to figure out the best way to approach this. The items are 1 ring, 3 fruits, and 1 chalk. My values are ring:6, fruit:1, chalk:1. The other player's values are unknown. The goal is to maximize the total score (mine plus theirs) if our proposals are complementary.} \\
\\ 
\texttt{First, I need to consider possible value functions the other player might have. Since the items are ring, fruits, and chalk, common valuations could be:} \\
\\ 
\texttt{1. They might value the ring highly (like me), maybe even more.} \\
\texttt{2. They might value fruits more, especially since there are 3 of them.} \\
\texttt{3. They might value chalk for some reason, maybe it's rare or has a special use.} \\
\\ 
\texttt{Since I don't know their values, I need to find a division that's efficient regardless. The key is to maximize the total value, so I should aim for a division where the sum of both players' values is highest.} \\
\\ 
\texttt{If I take the ring (6 for me), then the remaining items are 3 fruits and 1 chalk. If the other player values fruits at, say, 2 each, then giving them all 3 fruits would give them 6, and chalk maybe 1. But if they value fruits lower, maybe 1 each, then 3 fruits would be 3. But if they value chalk higher, maybe they prefer that.} \\
\\ 
\texttt{Alternatively, if the other player also wants the ring, there's a conflict. But since there's only one ring, we can't both have it. So maybe I should propose taking the ring and offer the other player the fruits and chalk. But if they need the ring, that's a problem.} \\
\\ 
\texttt{Wait, but the secret proposal needs to be complementary. So if I propose taking the ring, and they propose taking the fruits and chalk, that's okay. But if they also propose taking the ring, then our proposals conflict, and we both get zero.} \\
\\ 
\texttt{So maybe I should try to find out what they value. But since I can't ask directly (since messages are sent as is), maybe I can make a proposal that allows for different possibilities.} \\
\\ 
\texttt{Alternatively, maybe I can make a proposal that gives me the ring and the other player the rest. Let's see: my value would be 6 (ring) + 0 (since I don't take fruits or chalk). The other player's value would depend} \\
            }
        }
    }
     \\ \\

    \theutterance \stepcounter{utterance}  
    & & & \multicolumn{2}{p{0.3\linewidth}}{
        \cellcolor[rgb]{0.9,0.9,0.9}{
            \makecell[{{p{\linewidth}}}]{
                \texttt{\tiny{[GM$|$GM]}}
                \texttt{continue} \\
            }
        }
    }
    & & \\ \\

    \theutterance \stepcounter{utterance}  
    & & & \multicolumn{2}{p{0.3\linewidth}}{
        \cellcolor[rgb]{0.9,0.9,0.9}{
            \makecell[{{p{\linewidth}}}]{
                \texttt{\tiny{[GM$|$GM]}}
                \texttt{<think>} \\
\texttt{Okay, let's try to figure out the best way to approach this. The items are 1 ring, 3 fruits, and 1 chalk. My values are ring:6, fruit:1, chalk:1. The other player's values are unknown. The goal is to maximize the total score (mine plus theirs) if our proposals are complementary.} \\
\\ 
\texttt{First, I need to consider possible value functions the other player might have. Since the items are ring, fruits, and chalk, common valuations could be:} \\
\\ 
\texttt{1. They might value the ring highly (like me), maybe even more.} \\
\texttt{2. They might value fruits more, especially since there are 3 of them.} \\
\texttt{3. They might value chalk for some reason, maybe it's rare or has a special use.} \\
\\ 
\texttt{Since I don't know their values, I need to find a division that's efficient regardless. The key is to maximize the total value, so I should aim for a division where the sum of both players' values is highest.} \\
\\ 
\texttt{If I take the ring (6 for me), then the remaining items are 3 fruits and 1 chalk. If the other player values fruits at, say, 2 each, then giving them all 3 fruits would give them 6, and chalk maybe 1. But if they value fruits lower, maybe 1 each, then 3 fruits would be 3. But if they value chalk higher, maybe they prefer that.} \\
\\ 
\texttt{Alternatively, if the other player also wants the ring, there's a conflict. But since there's only one ring, we can't both have it. So maybe I should propose taking the ring and offer the other player the fruits and chalk. But if they need the ring, that's a problem.} \\
\\ 
\texttt{Wait, but the secret proposal needs to be complementary. So if I propose taking the ring, and they propose taking the fruits and chalk, that's okay. But if they also propose taking the ring, then our proposals conflict, and we both get zero.} \\
\\ 
\texttt{So maybe I should try to find out what they value. But since I can't ask directly (since messages are sent as is), maybe I can make a proposal that allows for different possibilities.} \\
\\ 
\texttt{Alternatively, maybe I can make a proposal that gives me the ring and the other player the rest. Let's see: my value would be 6 (ring) + 0 (since I don't take fruits or chalk). The other player's value would depend} \\
            }
        }
    }
    & & \\ \\

    \theutterance \stepcounter{utterance}  
    & & \multicolumn{4}{p{0.6\linewidth}}{
        \cellcolor[rgb]{0.9,0.9,0.9}{
            \makecell[{{p{\linewidth}}}]{
                \texttt{\tiny{[P1$\langle$GM]}}
                \texttt{<think>} \\
\texttt{Okay, let's try to figure out the best way to approach this. The items are 1 ring, 3 fruits, and 1 chalk. My values are ring:6, fruit:1, chalk:1. The other player's values are unknown. The goal is to maximize the total score (mine plus theirs) if our proposals are complementary.} \\
\\ 
\texttt{First, I need to consider possible value functions the other player might have. Since the items are ring, fruits, and chalk, common valuations could be:} \\
\\ 
\texttt{1. They might value the ring highly (like me), maybe even more.} \\
\texttt{2. They might value fruits more, especially since there are 3 of them.} \\
\texttt{3. They might value chalk for some reason, maybe it's rare or has a special use.} \\
\\ 
\texttt{Since I don't know their values, I need to find a division that's efficient regardless. The key is to maximize the total value, so I should aim for a division where the sum of both players' values is highest.} \\
\\ 
\texttt{If I take the ring (6 for me), then the remaining items are 3 fruits and 1 chalk. If the other player values fruits at, say, 2 each, then giving them all 3 fruits would give them 6, and chalk maybe 1. But if they value fruits lower, maybe 1 each, then 3 fruits would be 3. But if they value chalk higher, maybe they prefer that.} \\
\\ 
\texttt{Alternatively, if the other player also wants the ring, there's a conflict. But since there's only one ring, we can't both have it. So maybe I should propose taking the ring and offer the other player the fruits and chalk. But if they need the ring, that's a problem.} \\
\\ 
\texttt{Wait, but the secret proposal needs to be complementary. So if I propose taking the ring, and they propose taking the fruits and chalk, that's okay. But if they also propose taking the ring, then our proposals conflict, and we both get zero.} \\
\\ 
\texttt{So maybe I should try to find out what they value. But since I can't ask directly (since messages are sent as is), maybe I can make a proposal that allows for different possibilities.} \\
\\ 
\texttt{Alternatively, maybe I can make a proposal that gives me the ring and the other player the rest. Let's see: my value would be 6 (ring) + 0 (since I don't take fruits or chalk). The other player's value would depend} \\
            }
        }
    }
    & \\ \\

    \theutterance \stepcounter{utterance}  
    & \multicolumn{4}{p{0.6\linewidth}}{
        \cellcolor[rgb]{0.9,0.9,0.9}{
            \makecell[{{p{\linewidth}}}]{
                \texttt{\tiny{[P1$\rangle$GM]}}
                \texttt{<think>} \\
\texttt{Okay, let's try to figure out the best way to approach this. The items are 1 ring, 3 fruits, and 1 chalk. My values are ring:6, fruit:1, chalk:1. The other player's values are unknown. The goal is to maximize the total score (mine plus theirs) if our proposals are complementary.} \\
\\ 
\texttt{First, I need to consider possible value functions the other player might have. Since the items are ring, fruits, and chalk, common valuations could be:} \\
\\ 
\texttt{1. They might value the ring highly (like me), maybe even more.} \\
\texttt{2. They might value fruits more, especially since there are 3 of them.} \\
\texttt{3. They might value chalk for some reason, maybe it's rare or has a special use.} \\
\\ 
\texttt{Since I don't know their values, I need to find a division that's efficient regardless. The key is to maximize the total value, so I should aim for a division where the sum of both players' values is highest.} \\
\\ 
\texttt{If I take the ring (6 for me), then the remaining items are 3 fruits and 1 chalk. If the other player values fruits at, say, 2 each, then giving them all 3 fruits would give them 6, and chalk maybe 1. But if they value fruits lower, maybe 1 each, then 3 fruits would be 3. But if they value chalk higher, maybe they prefer that.} \\
\\ 
\texttt{Alternatively, if the other player also wants the ring, there's a conflict. But since there's only one ring, we can't both have it. So maybe I should propose taking the ring and offer the other player the fruits and chalk. But if they need the ring, that's a problem.} \\
\\ 
\texttt{Wait, but the secret proposal needs to be complementary. So if I propose taking the ring, and they propose taking the fruits and chalk, that's okay. But if they also propose taking the ring, then our proposals conflict, and we both get zero.} \\
\\ 
\texttt{So maybe I should try to find out what they value. But since I can't ask directly (since messages are sent as is), maybe I can make a proposal that allows for different possibilities.} \\
\\ 
\texttt{Alternatively, maybe I can make a proposal that gives me the ring and the other player the rest. Let's see: my value would be 6 (ring) + 0 (since I don't take fruits or chalk). The other player's value would depend} \\
            }
        }
    }
    & & \\ \\

    \theutterance \stepcounter{utterance}  
    & & & \multicolumn{2}{p{0.3\linewidth}}{
        \cellcolor[rgb]{0.9,0.9,0.9}{
            \makecell[{{p{\linewidth}}}]{
                \texttt{\tiny{[GM$|$GM]}}
                \texttt{continue} \\
            }
        }
    }
    & & \\ \\

    \theutterance \stepcounter{utterance}  
    & & & \multicolumn{2}{p{0.3\linewidth}}{
        \cellcolor[rgb]{0.9,0.9,0.9}{
            \makecell[{{p{\linewidth}}}]{
                \texttt{\tiny{[GM$|$GM]}}
                \texttt{<think>} \\
\texttt{Okay, let's try to figure out the best way to approach this. The items are 1 ring, 3 fruits, and 1 chalk. My values are ring:6, fruit:1, chalk:1. The other player's values are unknown. The goal is to maximize the total score (mine plus theirs) if our proposals are complementary.} \\
\\ 
\texttt{First, I need to consider possible value functions the other player might have. Since the items are ring, fruits, and chalk, common valuations could be:} \\
\\ 
\texttt{1. They might value the ring highly (like me), maybe even more.} \\
\texttt{2. They might value fruits more, especially since there are 3 of them.} \\
\texttt{3. They might value chalk for some reason, maybe it's rare or has a special use.} \\
\\ 
\texttt{Since I don't know their values, I need to find a division that's efficient regardless. The key is to maximize the total value, so I should aim for a division where the sum of both players' values is highest.} \\
\\ 
\texttt{If I take the ring (6 for me), then the remaining items are 3 fruits and 1 chalk. If the other player values fruits at, say, 2 each, then giving them all 3 fruits would give them 6, and chalk maybe 1. But if they value fruits lower, maybe 1 each, then 3 fruits would be 3. But if they value chalk higher, maybe they prefer that.} \\
\\ 
\texttt{Alternatively, if the other player also wants the ring, there's a conflict. But since there's only one ring, we can't both have it. So maybe I should propose taking the ring and offer the other player the fruits and chalk. But if they need the ring, that's a problem.} \\
\\ 
\texttt{Wait, but the secret proposal needs to be complementary. So if I propose taking the ring, and they propose taking the fruits and chalk, that's okay. But if they also propose taking the ring, then our proposals conflict, and we both get zero.} \\
\\ 
\texttt{So maybe I should try to find out what they value. But since I can't ask directly (since messages are sent as is), maybe I can make a proposal that allows for different possibilities.} \\
\\ 
\texttt{Alternatively, maybe I can make a proposal that gives me the ring and the other player the rest. Let's see: my value would be 6 (ring) + 0 (since I don't take fruits or chalk). The other player's value would depend} \\
            }
        }
    }
    & & \\ \\

    \theutterance \stepcounter{utterance}  
    & & \multicolumn{4}{p{0.6\linewidth}}{
        \cellcolor[rgb]{0.9,0.9,0.9}{
            \makecell[{{p{\linewidth}}}]{
                \texttt{\tiny{[GM$\rangle$P2]}}
                \texttt{<think>} \\
\texttt{Okay, let's try to figure out the best way to approach this. The items are 1 ring, 3 fruits, and 1 chalk. My values are ring:6, fruit:1, chalk:1. The other player's values are unknown. The goal is to maximize the total score (mine plus theirs) if our proposals are complementary.} \\
\\ 
\texttt{First, I need to consider possible value functions the other player might have. Since the items are ring, fruits, and chalk, common valuations could be:} \\
\\ 
\texttt{1. They might value the ring highly (like me), maybe even more.} \\
\texttt{2. They might value fruits more, especially since there are 3 of them.} \\
\texttt{3. They might value chalk for some reason, maybe it's rare or has a special use.} \\
\\ 
\texttt{Since I don't know their values, I need to find a division that's efficient regardless. The key is to maximize the total value, so I should aim for a division where the sum of both players' values is highest.} \\
\\ 
\texttt{If I take the ring (6 for me), then the remaining items are 3 fruits and 1 chalk. If the other player values fruits at, say, 2 each, then giving them all 3 fruits would give them 6, and chalk maybe 1. But if they value fruits lower, maybe 1 each, then 3 fruits would be 3. But if they value chalk higher, maybe they prefer that.} \\
\\ 
\texttt{Alternatively, if the other player also wants the ring, there's a conflict. But since there's only one ring, we can't both have it. So maybe I should propose taking the ring and offer the other player the fruits and chalk. But if they need the ring, that's a problem.} \\
\\ 
\texttt{Wait, but the secret proposal needs to be complementary. So if I propose taking the ring, and they propose taking the fruits and chalk, that's okay. But if they also propose taking the ring, then our proposals conflict, and we both get zero.} \\
\\ 
\texttt{So maybe I should try to find out what they value. But since I can't ask directly (since messages are sent as is), maybe I can make a proposal that allows for different possibilities.} \\
\\ 
\texttt{Alternatively, maybe I can make a proposal that gives me the ring and the other player the rest. Let's see: my value would be 6 (ring) + 0 (since I don't take fruits or chalk). The other player's value would depend} \\
            }
        }
    }
    & \\ \\

    \theutterance \stepcounter{utterance}  
    & & & \multicolumn{4}{p{0.6\linewidth}}{
        \cellcolor[rgb]{0.9,0.9,0.9}{
            \makecell[{{p{\linewidth}}}]{
                \texttt{\tiny{[GM$\langle$P2]}}
                \texttt{<think>} \\
\texttt{Okay, let's try to figure out the best way to approach this. The items are 1 ring, 3 fruits, and 1 chalk. My values are ring:6, fruit:1, chalk:1. The other player's values are unknown. The goal is to maximize the total score (mine plus theirs) if our proposals are complementary.} \\
\\ 
\texttt{First, I need to consider possible value functions the other player might have. Since the items are ring, fruits, and chalk, common valuations could be:} \\
\\ 
\texttt{1. They might value the ring highly (like me), maybe even more.} \\
\texttt{2. They might value fruits more, especially since there are 3 of them.} \\
\texttt{3. They might value chalk for some reason, maybe it's rare or has a special use.} \\
\\ 
\texttt{Since I don't know their values, I need to find a division that's efficient regardless. The key is to maximize the total value, so I should aim for a division where the sum of both players' values is highest.} \\
\\ 
\texttt{If I take the ring (6 for me), then the remaining items are 3 fruits and 1 chalk. If the other player values fruits at, say, 2 each, then giving them all 3 fruits would give them 6, and chalk maybe 1. But if they value fruits lower, maybe 1 each, then 3 fruits would be 3. But if they value chalk higher, maybe they prefer that.} \\
\\ 
\texttt{Alternatively, if the other player also wants the ring, there's a conflict. But since there's only one ring, we can't both have it. So maybe I should propose taking the ring and offer the other player the fruits and chalk. But if they need the ring, that's a problem.} \\
\\ 
\texttt{Wait, but the secret proposal needs to be complementary. So if I propose taking the ring, and they propose taking the fruits and chalk, that's okay. But if they also propose taking the ring, then our proposals conflict, and we both get zero.} \\
\\ 
\texttt{So maybe I should try to find out what they value. But since I can't ask directly (since messages are sent as is), maybe I can make a proposal that allows for different possibilities.} \\
\\ 
\texttt{Alternatively, maybe I can make a proposal that gives me the ring and the other player the rest. Let's see: my value would be 6 (ring) + 0 (since I don't take fruits or chalk). The other player's value would depend} \\
            }
        }
    }
     \\ \\

    \theutterance \stepcounter{utterance}  
    & & & \multicolumn{2}{p{0.3\linewidth}}{
        \cellcolor[rgb]{0.9,0.9,0.9}{
            \makecell[{{p{\linewidth}}}]{
                \texttt{\tiny{[GM$|$GM]}}
                \texttt{continue} \\
            }
        }
    }
    & & \\ \\

    \theutterance \stepcounter{utterance}  
    & & & \multicolumn{2}{p{0.3\linewidth}}{
        \cellcolor[rgb]{0.9,0.9,0.9}{
            \makecell[{{p{\linewidth}}}]{
                \texttt{\tiny{[GM$|$GM]}}
                \texttt{<think>} \\
\texttt{Okay, let's try to figure out the best way to approach this. The items are 1 ring, 3 fruits, and 1 chalk. My values are ring:6, fruit:1, chalk:1. The other player's values are unknown. The goal is to maximize the total score (mine plus theirs) if our proposals are complementary.} \\
\\ 
\texttt{First, I need to consider possible value functions the other player might have. Since the items are ring, fruits, and chalk, common valuations could be:} \\
\\ 
\texttt{1. They might value the ring highly (like me), maybe even more.} \\
\texttt{2. They might value fruits more, especially since there are 3 of them.} \\
\texttt{3. They might value chalk for some reason, maybe it's rare or has a special use.} \\
\\ 
\texttt{Since I don't know their values, I need to find a division that's efficient regardless. The key is to maximize the total value, so I should aim for a division where the sum of both players' values is highest.} \\
\\ 
\texttt{If I take the ring (6 for me), then the remaining items are 3 fruits and 1 chalk. If the other player values fruits at, say, 2 each, then giving them all 3 fruits would give them 6, and chalk maybe 1. But if they value fruits lower, maybe 1 each, then 3 fruits would be 3. But if they value chalk higher, maybe they prefer that.} \\
\\ 
\texttt{Alternatively, if the other player also wants the ring, there's a conflict. But since there's only one ring, we can't both have it. So maybe I should propose taking the ring and offer the other player the fruits and chalk. But if they need the ring, that's a problem.} \\
\\ 
\texttt{Wait, but the secret proposal needs to be complementary. So if I propose taking the ring, and they propose taking the fruits and chalk, that's okay. But if they also propose taking the ring, then our proposals conflict, and we both get zero.} \\
\\ 
\texttt{So maybe I should try to find out what they value. But since I can't ask directly (since messages are sent as is), maybe I can make a proposal that allows for different possibilities.} \\
\\ 
\texttt{Alternatively, maybe I can make a proposal that gives me the ring and the other player the rest. Let's see: my value would be 6 (ring) + 0 (since I don't take fruits or chalk). The other player's value would depend} \\
            }
        }
    }
    & & \\ \\

    \theutterance \stepcounter{utterance}  
    & & \multicolumn{4}{p{0.6\linewidth}}{
        \cellcolor[rgb]{0.9,0.9,0.9}{
            \makecell[{{p{\linewidth}}}]{
                \texttt{\tiny{[P1$\langle$GM]}}
                \texttt{<think>} \\
\texttt{Okay, let's try to figure out the best way to approach this. The items are 1 ring, 3 fruits, and 1 chalk. My values are ring:6, fruit:1, chalk:1. The other player's values are unknown. The goal is to maximize the total score (mine plus theirs) if our proposals are complementary.} \\
\\ 
\texttt{First, I need to consider possible value functions the other player might have. Since the items are ring, fruits, and chalk, common valuations could be:} \\
\\ 
\texttt{1. They might value the ring highly (like me), maybe even more.} \\
\texttt{2. They might value fruits more, especially since there are 3 of them.} \\
\texttt{3. They might value chalk for some reason, maybe it's rare or has a special use.} \\
\\ 
\texttt{Since I don't know their values, I need to find a division that's efficient regardless. The key is to maximize the total value, so I should aim for a division where the sum of both players' values is highest.} \\
\\ 
\texttt{If I take the ring (6 for me), then the remaining items are 3 fruits and 1 chalk. If the other player values fruits at, say, 2 each, then giving them all 3 fruits would give them 6, and chalk maybe 1. But if they value fruits lower, maybe 1 each, then 3 fruits would be 3. But if they value chalk higher, maybe they prefer that.} \\
\\ 
\texttt{Alternatively, if the other player also wants the ring, there's a conflict. But since there's only one ring, we can't both have it. So maybe I should propose taking the ring and offer the other player the fruits and chalk. But if they need the ring, that's a problem.} \\
\\ 
\texttt{Wait, but the secret proposal needs to be complementary. So if I propose taking the ring, and they propose taking the fruits and chalk, that's okay. But if they also propose taking the ring, then our proposals conflict, and we both get zero.} \\
\\ 
\texttt{So maybe I should try to find out what they value. But since I can't ask directly (since messages are sent as is), maybe I can make a proposal that allows for different possibilities.} \\
\\ 
\texttt{Alternatively, maybe I can make a proposal that gives me the ring and the other player the rest. Let's see: my value would be 6 (ring) + 0 (since I don't take fruits or chalk). The other player's value would depend} \\
            }
        }
    }
    & \\ \\

    \theutterance \stepcounter{utterance}  
    & \multicolumn{4}{p{0.6\linewidth}}{
        \cellcolor[rgb]{0.9,0.9,0.9}{
            \makecell[{{p{\linewidth}}}]{
                \texttt{\tiny{[P1$\rangle$GM]}}
                \texttt{<think>} \\
\texttt{Okay, let's try to figure out the best way to approach this. The items are 1 ring, 3 fruits, and 1 chalk. My values are ring:6, fruit:1, chalk:1. The other player's values are unknown. The goal is to maximize the total score (mine plus theirs) if our proposals are complementary.} \\
\\ 
\texttt{First, I need to consider possible value functions the other player might have. Since the items are ring, fruits, and chalk, common valuations could be:} \\
\\ 
\texttt{1. They might value the ring highly (like me), maybe even more.} \\
\texttt{2. They might value fruits more, especially since there are 3 of them.} \\
\texttt{3. They might value chalk for some reason, maybe it's rare or has a special use.} \\
\\ 
\texttt{Since I don't know their values, I need to find a division that's efficient regardless. The key is to maximize the total value, so I should aim for a division where the sum of both players' values is highest.} \\
\\ 
\texttt{If I take the ring (6 for me), then the remaining items are 3 fruits and 1 chalk. If the other player values fruits at, say, 2 each, then giving them all 3 fruits would give them 6, and chalk maybe 1. But if they value fruits lower, maybe 1 each, then 3 fruits would be 3. But if they value chalk higher, maybe they prefer that.} \\
\\ 
\texttt{Alternatively, if the other player also wants the ring, there's a conflict. But since there's only one ring, we can't both have it. So maybe I should propose taking the ring and offer the other player the fruits and chalk. But if they need the ring, that's a problem.} \\
\\ 
\texttt{Wait, but the secret proposal needs to be complementary. So if I propose taking the ring, and they propose taking the fruits and chalk, that's okay. But if they also propose taking the ring, then our proposals conflict, and we both get zero.} \\
\\ 
\texttt{So maybe I should try to find out what they value. But since I can't ask directly (since messages are sent as is), maybe I can make a proposal that allows for different possibilities.} \\
\\ 
\texttt{Alternatively, maybe I can make a proposal that gives me the ring and the other player the rest. Let's see: my value would be 6 (ring) + 0 (since I don't take fruits or chalk). The other player's value would depend} \\
            }
        }
    }
    & & \\ \\

    \theutterance \stepcounter{utterance}  
    & & & \multicolumn{2}{p{0.3\linewidth}}{
        \cellcolor[rgb]{0.9,0.9,0.9}{
            \makecell[{{p{\linewidth}}}]{
                \texttt{\tiny{[GM$|$GM]}}
                \texttt{continue} \\
            }
        }
    }
    & & \\ \\

    \theutterance \stepcounter{utterance}  
    & & & \multicolumn{2}{p{0.3\linewidth}}{
        \cellcolor[rgb]{0.9,0.9,0.9}{
            \makecell[{{p{\linewidth}}}]{
                \texttt{\tiny{[GM$|$GM]}}
                \texttt{<think>} \\
\texttt{Okay, let's try to figure out the best way to approach this. The items are 1 ring, 3 fruits, and 1 chalk. My values are ring:6, fruit:1, chalk:1. The other player's values are unknown. The goal is to maximize the total score (mine plus theirs) if our proposals are complementary.} \\
\\ 
\texttt{First, I need to consider possible value functions the other player might have. Since the items are ring, fruits, and chalk, common valuations could be:} \\
\\ 
\texttt{1. They might value the ring highly (like me), maybe even more.} \\
\texttt{2. They might value fruits more, especially since there are 3 of them.} \\
\texttt{3. They might value chalk for some reason, maybe it's rare or has a special use.} \\
\\ 
\texttt{Since I don't know their values, I need to find a division that's efficient regardless. The key is to maximize the total value, so I should aim for a division where the sum of both players' values is highest.} \\
\\ 
\texttt{If I take the ring (6 for me), then the remaining items are 3 fruits and 1 chalk. If the other player values fruits at, say, 2 each, then giving them all 3 fruits would give them 6, and chalk maybe 1. But if they value fruits lower, maybe 1 each, then 3 fruits would be 3. But if they value chalk higher, maybe they prefer that.} \\
\\ 
\texttt{Alternatively, if the other player also wants the ring, there's a conflict. But since there's only one ring, we can't both have it. So maybe I should propose taking the ring and offer the other player the fruits and chalk. But if they need the ring, that's a problem.} \\
\\ 
\texttt{Wait, but the secret proposal needs to be complementary. So if I propose taking the ring, and they propose taking the fruits and chalk, that's okay. But if they also propose taking the ring, then our proposals conflict, and we both get zero.} \\
\\ 
\texttt{So maybe I should try to find out what they value. But since I can't ask directly (since messages are sent as is), maybe I can make a proposal that allows for different possibilities.} \\
\\ 
\texttt{Alternatively, maybe I can make a proposal that gives me the ring and the other player the rest. Let's see: my value would be 6 (ring) + 0 (since I don't take fruits or chalk). The other player's value would depend} \\
            }
        }
    }
    & & \\ \\

    \theutterance \stepcounter{utterance}  
    & & \multicolumn{4}{p{0.6\linewidth}}{
        \cellcolor[rgb]{0.9,0.9,0.9}{
            \makecell[{{p{\linewidth}}}]{
                \texttt{\tiny{[GM$\rangle$P2]}}
                \texttt{<think>} \\
\texttt{Okay, let's try to figure out the best way to approach this. The items are 1 ring, 3 fruits, and 1 chalk. My values are ring:6, fruit:1, chalk:1. The other player's values are unknown. The goal is to maximize the total score (mine plus theirs) if our proposals are complementary.} \\
\\ 
\texttt{First, I need to consider possible value functions the other player might have. Since the items are ring, fruits, and chalk, common valuations could be:} \\
\\ 
\texttt{1. They might value the ring highly (like me), maybe even more.} \\
\texttt{2. They might value fruits more, especially since there are 3 of them.} \\
\texttt{3. They might value chalk for some reason, maybe it's rare or has a special use.} \\
\\ 
\texttt{Since I don't know their values, I need to find a division that's efficient regardless. The key is to maximize the total value, so I should aim for a division where the sum of both players' values is highest.} \\
\\ 
\texttt{If I take the ring (6 for me), then the remaining items are 3 fruits and 1 chalk. If the other player values fruits at, say, 2 each, then giving them all 3 fruits would give them 6, and chalk maybe 1. But if they value fruits lower, maybe 1 each, then 3 fruits would be 3. But if they value chalk higher, maybe they prefer that.} \\
\\ 
\texttt{Alternatively, if the other player also wants the ring, there's a conflict. But since there's only one ring, we can't both have it. So maybe I should propose taking the ring and offer the other player the fruits and chalk. But if they need the ring, that's a problem.} \\
\\ 
\texttt{Wait, but the secret proposal needs to be complementary. So if I propose taking the ring, and they propose taking the fruits and chalk, that's okay. But if they also propose taking the ring, then our proposals conflict, and we both get zero.} \\
\\ 
\texttt{So maybe I should try to find out what they value. But since I can't ask directly (since messages are sent as is), maybe I can make a proposal that allows for different possibilities.} \\
\\ 
\texttt{Alternatively, maybe I can make a proposal that gives me the ring and the other player the rest. Let's see: my value would be 6 (ring) + 0 (since I don't take fruits or chalk). The other player's value would depend} \\
            }
        }
    }
    & \\ \\

    \theutterance \stepcounter{utterance}  
    & & & \multicolumn{4}{p{0.6\linewidth}}{
        \cellcolor[rgb]{0.9,0.9,0.9}{
            \makecell[{{p{\linewidth}}}]{
                \texttt{\tiny{[GM$\langle$P2]}}
                \texttt{<think>} \\
\texttt{Okay, let's try to figure out the best way to approach this. The items are 1 ring, 3 fruits, and 1 chalk. My values are ring:6, fruit:1, chalk:1. The other player's values are unknown. The goal is to maximize the total score (mine plus theirs) if our proposals are complementary.} \\
\\ 
\texttt{First, I need to consider possible value functions the other player might have. Since the items are ring, fruits, and chalk, common valuations could be:} \\
\\ 
\texttt{1. They might value the ring highly (like me), maybe even more.} \\
\texttt{2. They might value fruits more, especially since there are 3 of them.} \\
\texttt{3. They might value chalk for some reason, maybe it's rare or has a special use.} \\
\\ 
\texttt{Since I don't know their values, I need to find a division that's efficient regardless. The key is to maximize the total value, so I should aim for a division where the sum of both players' values is highest.} \\
\\ 
\texttt{If I take the ring (6 for me), then the remaining items are 3 fruits and 1 chalk. If the other player values fruits at, say, 2 each, then giving them all 3 fruits would give them 6, and chalk maybe 1. But if they value fruits lower, maybe 1 each, then 3 fruits would be 3. But if they value chalk higher, maybe they prefer that.} \\
\\ 
\texttt{Alternatively, if the other player also wants the ring, there's a conflict. But since there's only one ring, we can't both have it. So maybe I should propose taking the ring and offer the other player the fruits and chalk. But if they need the ring, that's a problem.} \\
\\ 
\texttt{Wait, but the secret proposal needs to be complementary. So if I propose taking the ring, and they propose taking the fruits and chalk, that's okay. But if they also propose taking the ring, then our proposals conflict, and we both get zero.} \\
\\ 
\texttt{So maybe I should try to find out what they value. But since I can't ask directly (since messages are sent as is), maybe I can make a proposal that allows for different possibilities.} \\
\\ 
\texttt{Alternatively, maybe I can make a proposal that gives me the ring and the other player the rest. Let's see: my value would be 6 (ring) + 0 (since I don't take fruits or chalk). The other player's value would depend} \\
            }
        }
    }
     \\ \\

    \theutterance \stepcounter{utterance}  
    & & & \multicolumn{2}{p{0.3\linewidth}}{
        \cellcolor[rgb]{0.9,0.9,0.9}{
            \makecell[{{p{\linewidth}}}]{
                \texttt{\tiny{[GM$|$GM]}}
                \texttt{continue} \\
            }
        }
    }
    & & \\ \\

    \theutterance \stepcounter{utterance}  
    & & & \multicolumn{2}{p{0.3\linewidth}}{
        \cellcolor[rgb]{0.9,0.9,0.9}{
            \makecell[{{p{\linewidth}}}]{
                \texttt{\tiny{[GM$|$GM]}}
                \texttt{<think>} \\
\texttt{Okay, let's try to figure out the best way to approach this. The items are 1 ring, 3 fruits, and 1 chalk. My values are ring:6, fruit:1, chalk:1. The other player's values are unknown. The goal is to maximize the total score (mine plus theirs) if our proposals are complementary.} \\
\\ 
\texttt{First, I need to consider possible value functions the other player might have. Since the items are ring, fruits, and chalk, common valuations could be:} \\
\\ 
\texttt{1. They might value the ring highly (like me), maybe even more.} \\
\texttt{2. They might value fruits more, especially since there are 3 of them.} \\
\texttt{3. They might value chalk for some reason, maybe it's rare or has a special use.} \\
\\ 
\texttt{Since I don't know their values, I need to find a division that's efficient regardless. The key is to maximize the total value, so I should aim for a division where the sum of both players' values is highest.} \\
\\ 
\texttt{If I take the ring (6 for me), then the remaining items are 3 fruits and 1 chalk. If the other player values fruits at, say, 2 each, then giving them all 3 fruits would give them 6, and chalk maybe 1. But if they value fruits lower, maybe 1 each, then 3 fruits would be 3. But if they value chalk higher, maybe they prefer that.} \\
\\ 
\texttt{Alternatively, if the other player also wants the ring, there's a conflict. But since there's only one ring, we can't both have it. So maybe I should propose taking the ring and offer the other player the fruits and chalk. But if they need the ring, that's a problem.} \\
\\ 
\texttt{Wait, but the secret proposal needs to be complementary. So if I propose taking the ring, and they propose taking the fruits and chalk, that's okay. But if they also propose taking the ring, then our proposals conflict, and we both get zero.} \\
\\ 
\texttt{So maybe I should try to find out what they value. But since I can't ask directly (since messages are sent as is), maybe I can make a proposal that allows for different possibilities.} \\
\\ 
\texttt{Alternatively, maybe I can make a proposal that gives me the ring and the other player the rest. Let's see: my value would be 6 (ring) + 0 (since I don't take fruits or chalk). The other player's value would depend} \\
            }
        }
    }
    & & \\ \\

    \theutterance \stepcounter{utterance}  
    & & \multicolumn{4}{p{0.6\linewidth}}{
        \cellcolor[rgb]{0.9,0.9,0.9}{
            \makecell[{{p{\linewidth}}}]{
                \texttt{\tiny{[P1$\langle$GM]}}
                \texttt{<think>} \\
\texttt{Okay, let's try to figure out the best way to approach this. The items are 1 ring, 3 fruits, and 1 chalk. My values are ring:6, fruit:1, chalk:1. The other player's values are unknown. The goal is to maximize the total score (mine plus theirs) if our proposals are complementary.} \\
\\ 
\texttt{First, I need to consider possible value functions the other player might have. Since the items are ring, fruits, and chalk, common valuations could be:} \\
\\ 
\texttt{1. They might value the ring highly (like me), maybe even more.} \\
\texttt{2. They might value fruits more, especially since there are 3 of them.} \\
\texttt{3. They might value chalk for some reason, maybe it's rare or has a special use.} \\
\\ 
\texttt{Since I don't know their values, I need to find a division that's efficient regardless. The key is to maximize the total value, so I should aim for a division where the sum of both players' values is highest.} \\
\\ 
\texttt{If I take the ring (6 for me), then the remaining items are 3 fruits and 1 chalk. If the other player values fruits at, say, 2 each, then giving them all 3 fruits would give them 6, and chalk maybe 1. But if they value fruits lower, maybe 1 each, then 3 fruits would be 3. But if they value chalk higher, maybe they prefer that.} \\
\\ 
\texttt{Alternatively, if the other player also wants the ring, there's a conflict. But since there's only one ring, we can't both have it. So maybe I should propose taking the ring and offer the other player the fruits and chalk. But if they need the ring, that's a problem.} \\
\\ 
\texttt{Wait, but the secret proposal needs to be complementary. So if I propose taking the ring, and they propose taking the fruits and chalk, that's okay. But if they also propose taking the ring, then our proposals conflict, and we both get zero.} \\
\\ 
\texttt{So maybe I should try to find out what they value. But since I can't ask directly (since messages are sent as is), maybe I can make a proposal that allows for different possibilities.} \\
\\ 
\texttt{Alternatively, maybe I can make a proposal that gives me the ring and the other player the rest. Let's see: my value would be 6 (ring) + 0 (since I don't take fruits or chalk). The other player's value would depend} \\
            }
        }
    }
    & \\ \\

    \theutterance \stepcounter{utterance}  
    & \multicolumn{4}{p{0.6\linewidth}}{
        \cellcolor[rgb]{0.9,0.9,0.9}{
            \makecell[{{p{\linewidth}}}]{
                \texttt{\tiny{[P1$\rangle$GM]}}
                \texttt{<think>} \\
\texttt{Okay, let's try to figure out the best way to approach this. The items are 1 ring, 3 fruits, and 1 chalk. My values are ring:6, fruit:1, chalk:1. The other player's values are unknown. The goal is to maximize the total score (mine plus theirs) if our proposals are complementary.} \\
\\ 
\texttt{First, I need to consider possible value functions the other player might have. Since the items are ring, fruits, and chalk, common valuations could be:} \\
\\ 
\texttt{1. They might value the ring highly (like me), maybe even more.} \\
\texttt{2. They might value fruits more, especially since there are 3 of them.} \\
\texttt{3. They might value chalk for some reason, maybe it's rare or has a special use.} \\
\\ 
\texttt{Since I don't know their values, I need to find a division that's efficient regardless. The key is to maximize the total value, so I should aim for a division where the sum of both players' values is highest.} \\
\\ 
\texttt{If I take the ring (6 for me), then the remaining items are 3 fruits and 1 chalk. If the other player values fruits at, say, 2 each, then giving them all 3 fruits would give them 6, and chalk maybe 1. But if they value fruits lower, maybe 1 each, then 3 fruits would be 3. But if they value chalk higher, maybe they prefer that.} \\
\\ 
\texttt{Alternatively, if the other player also wants the ring, there's a conflict. But since there's only one ring, we can't both have it. So maybe I should propose taking the ring and offer the other player the fruits and chalk. But if they need the ring, that's a problem.} \\
\\ 
\texttt{Wait, but the secret proposal needs to be complementary. So if I propose taking the ring, and they propose taking the fruits and chalk, that's okay. But if they also propose taking the ring, then our proposals conflict, and we both get zero.} \\
\\ 
\texttt{So maybe I should try to find out what they value. But since I can't ask directly (since messages are sent as is), maybe I can make a proposal that allows for different possibilities.} \\
\\ 
\texttt{Alternatively, maybe I can make a proposal that gives me the ring and the other player the rest. Let's see: my value would be 6 (ring) + 0 (since I don't take fruits or chalk). The other player's value would depend} \\
            }
        }
    }
    & & \\ \\

    \theutterance \stepcounter{utterance}  
    & & & \multicolumn{2}{p{0.3\linewidth}}{
        \cellcolor[rgb]{0.9,0.9,0.9}{
            \makecell[{{p{\linewidth}}}]{
                \texttt{\tiny{[GM$|$GM]}}
                \texttt{continue} \\
            }
        }
    }
    & & \\ \\

    \theutterance \stepcounter{utterance}  
    & & & \multicolumn{2}{p{0.3\linewidth}}{
        \cellcolor[rgb]{0.9,0.9,0.9}{
            \makecell[{{p{\linewidth}}}]{
                \texttt{\tiny{[GM$|$GM]}}
                \texttt{<think>} \\
\texttt{Okay, let's try to figure out the best way to approach this. The items are 1 ring, 3 fruits, and 1 chalk. My values are ring:6, fruit:1, chalk:1. The other player's values are unknown. The goal is to maximize the total score (mine plus theirs) if our proposals are complementary.} \\
\\ 
\texttt{First, I need to consider possible value functions the other player might have. Since the items are ring, fruits, and chalk, common valuations could be:} \\
\\ 
\texttt{1. They might value the ring highly (like me), maybe even more.} \\
\texttt{2. They might value fruits more, especially since there are 3 of them.} \\
\texttt{3. They might value chalk for some reason, maybe it's rare or has a special use.} \\
\\ 
\texttt{Since I don't know their values, I need to find a division that's efficient regardless. The key is to maximize the total value, so I should aim for a division where the sum of both players' values is highest.} \\
\\ 
\texttt{If I take the ring (6 for me), then the remaining items are 3 fruits and 1 chalk. If the other player values fruits at, say, 2 each, then giving them all 3 fruits would give them 6, and chalk maybe 1. But if they value fruits lower, maybe 1 each, then 3 fruits would be 3. But if they value chalk higher, maybe they prefer that.} \\
\\ 
\texttt{Alternatively, if the other player also wants the ring, there's a conflict. But since there's only one ring, we can't both have it. So maybe I should propose taking the ring and offer the other player the fruits and chalk. But if they need the ring, that's a problem.} \\
\\ 
\texttt{Wait, but the secret proposal needs to be complementary. So if I propose taking the ring, and they propose taking the fruits and chalk, that's okay. But if they also propose taking the ring, then our proposals conflict, and we both get zero.} \\
\\ 
\texttt{So maybe I should try to find out what they value. But since I can't ask directly (since messages are sent as is), maybe I can make a proposal that allows for different possibilities.} \\
\\ 
\texttt{Alternatively, maybe I can make a proposal that gives me the ring and the other player the rest. Let's see: my value would be 6 (ring) + 0 (since I don't take fruits or chalk). The other player's value would depend} \\
            }
        }
    }
    & & \\ \\

    \theutterance \stepcounter{utterance}  
    & & \multicolumn{4}{p{0.6\linewidth}}{
        \cellcolor[rgb]{0.9,0.9,0.9}{
            \makecell[{{p{\linewidth}}}]{
                \texttt{\tiny{[GM$\rangle$P2]}}
                \texttt{<think>} \\
\texttt{Okay, let's try to figure out the best way to approach this. The items are 1 ring, 3 fruits, and 1 chalk. My values are ring:6, fruit:1, chalk:1. The other player's values are unknown. The goal is to maximize the total score (mine plus theirs) if our proposals are complementary.} \\
\\ 
\texttt{First, I need to consider possible value functions the other player might have. Since the items are ring, fruits, and chalk, common valuations could be:} \\
\\ 
\texttt{1. They might value the ring highly (like me), maybe even more.} \\
\texttt{2. They might value fruits more, especially since there are 3 of them.} \\
\texttt{3. They might value chalk for some reason, maybe it's rare or has a special use.} \\
\\ 
\texttt{Since I don't know their values, I need to find a division that's efficient regardless. The key is to maximize the total value, so I should aim for a division where the sum of both players' values is highest.} \\
\\ 
\texttt{If I take the ring (6 for me), then the remaining items are 3 fruits and 1 chalk. If the other player values fruits at, say, 2 each, then giving them all 3 fruits would give them 6, and chalk maybe 1. But if they value fruits lower, maybe 1 each, then 3 fruits would be 3. But if they value chalk higher, maybe they prefer that.} \\
\\ 
\texttt{Alternatively, if the other player also wants the ring, there's a conflict. But since there's only one ring, we can't both have it. So maybe I should propose taking the ring and offer the other player the fruits and chalk. But if they need the ring, that's a problem.} \\
\\ 
\texttt{Wait, but the secret proposal needs to be complementary. So if I propose taking the ring, and they propose taking the fruits and chalk, that's okay. But if they also propose taking the ring, then our proposals conflict, and we both get zero.} \\
\\ 
\texttt{So maybe I should try to find out what they value. But since I can't ask directly (since messages are sent as is), maybe I can make a proposal that allows for different possibilities.} \\
\\ 
\texttt{Alternatively, maybe I can make a proposal that gives me the ring and the other player the rest. Let's see: my value would be 6 (ring) + 0 (since I don't take fruits or chalk). The other player's value would depend} \\
            }
        }
    }
    & \\ \\

    \theutterance \stepcounter{utterance}  
    & & & \multicolumn{4}{p{0.6\linewidth}}{
        \cellcolor[rgb]{0.9,0.9,0.9}{
            \makecell[{{p{\linewidth}}}]{
                \texttt{\tiny{[GM$\langle$P2]}}
                \texttt{<think>} \\
\texttt{Okay, let's try to figure out the best way to approach this. The items are 1 ring, 3 fruits, and 1 chalk. My values are ring:6, fruit:1, chalk:1. The other player's values are unknown. The goal is to maximize the total score (mine plus theirs) if our proposals are complementary.} \\
\\ 
\texttt{First, I need to consider possible value functions the other player might have. Since the items are ring, fruits, and chalk, common valuations could be:} \\
\\ 
\texttt{1. They might value the ring highly (like me), maybe even more.} \\
\texttt{2. They might value fruits more, especially since there are 3 of them.} \\
\texttt{3. They might value chalk for some reason, maybe it's rare or has a special use.} \\
\\ 
\texttt{Since I don't know their values, I need to find a division that's efficient regardless. The key is to maximize the total value, so I should aim for a division where the sum of both players' values is highest.} \\
\\ 
\texttt{If I take the ring (6 for me), then the remaining items are 3 fruits and 1 chalk. If the other player values fruits at, say, 2 each, then giving them all 3 fruits would give them 6, and chalk maybe 1. But if they value fruits lower, maybe 1 each, then 3 fruits would be 3. But if they value chalk higher, maybe they prefer that.} \\
\\ 
\texttt{Alternatively, if the other player also wants the ring, there's a conflict. But since there's only one ring, we can't both have it. So maybe I should propose taking the ring and offer the other player the fruits and chalk. But if they need the ring, that's a problem.} \\
\\ 
\texttt{Wait, but the secret proposal needs to be complementary. So if I propose taking the ring, and they propose taking the fruits and chalk, that's okay. But if they also propose taking the ring, then our proposals conflict, and we both get zero.} \\
\\ 
\texttt{So maybe I should try to find out what they value. But since I can't ask directly (since messages are sent as is), maybe I can make a proposal that allows for different possibilities.} \\
\\ 
\texttt{Alternatively, maybe I can make a proposal that gives me the ring and the other player the rest. Let's see: my value would be 6 (ring) + 0 (since I don't take fruits or chalk). The other player's value would depend} \\
            }
        }
    }
     \\ \\

    \theutterance \stepcounter{utterance}  
    & & & \multicolumn{2}{p{0.3\linewidth}}{
        \cellcolor[rgb]{0.9,0.9,0.9}{
            \makecell[{{p{\linewidth}}}]{
                \texttt{\tiny{[GM$|$GM]}}
                \texttt{continue} \\
            }
        }
    }
    & & \\ \\

    \theutterance \stepcounter{utterance}  
    & & & \multicolumn{2}{p{0.3\linewidth}}{
        \cellcolor[rgb]{0.9,0.9,0.9}{
            \makecell[{{p{\linewidth}}}]{
                \texttt{\tiny{[GM$|$GM]}}
                \texttt{<think>} \\
\texttt{Okay, let's try to figure out the best way to approach this. The items are 1 ring, 3 fruits, and 1 chalk. My values are ring:6, fruit:1, chalk:1. The other player's values are unknown. The goal is to maximize the total score (mine plus theirs) if our proposals are complementary.} \\
\\ 
\texttt{First, I need to consider possible value functions the other player might have. Since the items are ring, fruits, and chalk, common valuations could be:} \\
\\ 
\texttt{1. They might value the ring highly (like me), maybe even more.} \\
\texttt{2. They might value fruits more, especially since there are 3 of them.} \\
\texttt{3. They might value chalk for some reason, maybe it's rare or has a special use.} \\
\\ 
\texttt{Since I don't know their values, I need to find a division that's efficient regardless. The key is to maximize the total value, so I should aim for a division where the sum of both players' values is highest.} \\
\\ 
\texttt{If I take the ring (6 for me), then the remaining items are 3 fruits and 1 chalk. If the other player values fruits at, say, 2 each, then giving them all 3 fruits would give them 6, and chalk maybe 1. But if they value fruits lower, maybe 1 each, then 3 fruits would be 3. But if they value chalk higher, maybe they prefer that.} \\
\\ 
\texttt{Alternatively, if the other player also wants the ring, there's a conflict. But since there's only one ring, we can't both have it. So maybe I should propose taking the ring and offer the other player the fruits and chalk. But if they need the ring, that's a problem.} \\
\\ 
\texttt{Wait, but the secret proposal needs to be complementary. So if I propose taking the ring, and they propose taking the fruits and chalk, that's okay. But if they also propose taking the ring, then our proposals conflict, and we both get zero.} \\
\\ 
\texttt{So maybe I should try to find out what they value. But since I can't ask directly (since messages are sent as is), maybe I can make a proposal that allows for different possibilities.} \\
\\ 
\texttt{Alternatively, maybe I can make a proposal that gives me the ring and the other player the rest. Let's see: my value would be 6 (ring) + 0 (since I don't take fruits or chalk). The other player's value would depend} \\
            }
        }
    }
    & & \\ \\

    \theutterance \stepcounter{utterance}  
    & & \multicolumn{4}{p{0.6\linewidth}}{
        \cellcolor[rgb]{0.9,0.9,0.9}{
            \makecell[{{p{\linewidth}}}]{
                \texttt{\tiny{[P1$\langle$GM]}}
                \texttt{<think>} \\
\texttt{Okay, let's try to figure out the best way to approach this. The items are 1 ring, 3 fruits, and 1 chalk. My values are ring:6, fruit:1, chalk:1. The other player's values are unknown. The goal is to maximize the total score (mine plus theirs) if our proposals are complementary.} \\
\\ 
\texttt{First, I need to consider possible value functions the other player might have. Since the items are ring, fruits, and chalk, common valuations could be:} \\
\\ 
\texttt{1. They might value the ring highly (like me), maybe even more.} \\
\texttt{2. They might value fruits more, especially since there are 3 of them.} \\
\texttt{3. They might value chalk for some reason, maybe it's rare or has a special use.} \\
\\ 
\texttt{Since I don't know their values, I need to find a division that's efficient regardless. The key is to maximize the total value, so I should aim for a division where the sum of both players' values is highest.} \\
\\ 
\texttt{If I take the ring (6 for me), then the remaining items are 3 fruits and 1 chalk. If the other player values fruits at, say, 2 each, then giving them all 3 fruits would give them 6, and chalk maybe 1. But if they value fruits lower, maybe 1 each, then 3 fruits would be 3. But if they value chalk higher, maybe they prefer that.} \\
\\ 
\texttt{Alternatively, if the other player also wants the ring, there's a conflict. But since there's only one ring, we can't both have it. So maybe I should propose taking the ring and offer the other player the fruits and chalk. But if they need the ring, that's a problem.} \\
\\ 
\texttt{Wait, but the secret proposal needs to be complementary. So if I propose taking the ring, and they propose taking the fruits and chalk, that's okay. But if they also propose taking the ring, then our proposals conflict, and we both get zero.} \\
\\ 
\texttt{So maybe I should try to find out what they value. But since I can't ask directly (since messages are sent as is), maybe I can make a proposal that allows for different possibilities.} \\
\\ 
\texttt{Alternatively, maybe I can make a proposal that gives me the ring and the other player the rest. Let's see: my value would be 6 (ring) + 0 (since I don't take fruits or chalk). The other player's value would depend} \\
\\ 
\\ 
\texttt{The time is up. It is now your turn to submit a secret proposal.} \\
            }
        }
    }
    & \\ \\

    \theutterance \stepcounter{utterance}  
    & \multicolumn{4}{p{0.6\linewidth}}{
        \cellcolor[rgb]{0.9,0.9,0.9}{
            \makecell[{{p{\linewidth}}}]{
                \texttt{\tiny{[P1$\rangle$GM]}}
                \texttt{<think>} \\
\texttt{Okay, let's try to figure out the best way to approach this. The items are 1 ring, 3 fruits, and 1 chalk. My values are ring:6, fruit:1, chalk:1. The other player's values are unknown. The goal is to maximize the total score (mine plus theirs) if our proposals are complementary.} \\
\\ 
\texttt{First, I need to consider possible value functions the other player might have. Since the items are ring, fruits, and chalk, common valuations could be:} \\
\\ 
\texttt{1. They might value the ring highly (like me), maybe even more.} \\
\texttt{2. They might value fruits more, especially since there are 3 of them.} \\
\texttt{3. They might value chalk for some reason, maybe it's rare or has a special use.} \\
\\ 
\texttt{Since I don't know their values, I need to find a division that's efficient regardless. The key is to maximize the total value, so I should aim for a division where the sum of both players' values is highest.} \\
\\ 
\texttt{If I take the ring (6 for me), then the remaining items are 3 fruits and 1 chalk. If the other player values fruits at, say, 2 each, then giving them all 3 fruits would give them 6, and chalk maybe 1. But if they value fruits lower, maybe 1 each, then 3 fruits would be 3. But if they value chalk higher, maybe they prefer that.} \\
\\ 
\texttt{Alternatively, if the other player also wants the ring, there's a conflict. But since there's only one ring, we can't both have it. So maybe I should propose taking the ring and offer the other player the fruits and chalk. But if they need the ring, that's a problem.} \\
\\ 
\texttt{Wait, but the secret proposal needs to be complementary. So if I propose taking the ring, and they propose taking the fruits and chalk, that's okay. But if they also propose taking the ring, then our proposals conflict, and we both get zero.} \\
\\ 
\texttt{So maybe I should try to find out what they value. But since I can't ask directly (since messages are sent as is), maybe I can make a proposal that allows for different possibilities.} \\
\\ 
\texttt{Alternatively, maybe I can make a proposal that gives me the ring and the other player the rest. Let's see: my value would be 6 (ring) + 0 (since I don't take fruits or chalk). The other player's value would depend} \\
            }
        }
    }
    & & \\ \\

    \theutterance \stepcounter{utterance}  
    & & & \multicolumn{2}{p{0.3\linewidth}}{
        \cellcolor[rgb]{0.9,0.9,0.9}{
            \makecell[{{p{\linewidth}}}]{
                \texttt{\tiny{[GM$|$GM]}}
                \texttt{continue} \\
            }
        }
    }
    & & \\ \\

    \theutterance \stepcounter{utterance}  
    & & & \multicolumn{2}{p{0.3\linewidth}}{
        \cellcolor[rgb]{0.9,0.9,0.9}{
            \makecell[{{p{\linewidth}}}]{
                \texttt{\tiny{[GM$|$GM]}}
                \texttt{player was instructed to make a proposal} \\
            }
        }
    }
    & & \\ \\

\end{supertabular}
}

\end{document}
